% Generated draft: requires mathematical and source review.
\documentclass[11pt]{article}
\usepackage[margin=25mm]{geometry}
\usepackage{fontspec}
\setmainfont{Latin Modern Roman}
\usepackage{amsmath,amssymb,mathtools,unicode-math}
\setmathfont{Latin Modern Math}
\usepackage{enumitem,needspace,xurl,hyperref}
\setlength{\parindent}{0pt}
\setlength{\parskip}{5pt}
\setlength{\emergencystretch}{3em}
\newcommand{\N}{\mathbb{N}}
\newcommand{\Z}{\mathbb{Z}}
\newcommand{\Q}{\mathbb{Q}}
\newcommand{\R}{\mathbb{R}}
\newcommand{\C}{\mathbb{C}}
\newcommand{\F}{\mathbb{F}}
\newcommand{\A}{\mathbb{A}}
\newcommand{\PP}{\mathbb{P}}
\newcommand{\E}{\mathbb{E}}
\newcommand{\eps}{\varepsilon}
\newcommand{\dd}{\,\mathrm{d}}
\begin{document}
\section*{UnsolvedMath ID: 8800070}
\textbf{Source problem code:} catalogue entry
\subsection{Definitions and mathematical statement}
An ideal lattice is the image of a nonzero ideal in a specified order of a number field under a specified real coordinate embedding; the field family, order and embedding are part of any algorithmic instance definition.

A cryptographic pseudorandom generator is a uniformly efficiently computable family
\[
 G_\lambda:\{0,1\}^{s(\lambda)}\to\{0,1\}^{m(\lambda)},\qquad m(\lambda)>s(\lambda),
\]
such that for every probabilistic polynomial-time distinguisher \(D\),
\[
 \left|\Pr[D(1^\lambda,G_\lambda(U_{s(\lambda)}))=1]
 -\Pr[D(1^\lambda,U_{m(\lambda)})=1]\right|
\]
is negligible in the security parameter \(\lambda\). Here \(U_t\) is uniform on \(t\)-bit strings, and negligible means smaller than every inverse positive polynomial for all sufficiently large \(\lambda\).

Micciancio asks to construct efficient pseudorandom generators from ideal-lattice problems: specify the generator and prove its pseudorandomness by a reduction to an explicit computational hardness assumption on a specified ideal-lattice family. The statement does not identify the underlying search or decision problem, norm, approximation factor, parameter distribution, stretch, concrete efficiency target, or classical versus quantum reduction model.

The registry reports the construction goal solved through cyclic-lattice work and ring-LWE pseudorandomness. These are conditional cryptographic security claims, not proofs that ideal-lattice problems are unconditionally hard. Pseudorandom ring-LWE samples alone do not specify the seed accounting, expansion, and parameter choices of the requested bit-string generator; the concrete construction and security reduction need the cited sources.
\subsection{Short English statement}
Construct expanding bit-string generators whose efficient indistinguishability from uniform follows from a precisely stated ideal-lattice hardness assumption. The supplied resolution is reported through cyclic lattices and ring-LWE, with construction and reduction details still needing verification.
\subsection{Sources}
\begin{itemize}
\item Local UnsolvedMath registry, permanent ID 8800070, source code catalogue entry; attributed to Daniele Micciancio; no proposal year supplied; update 2026-08-06; dated status solved.
\item IPAM, \emph{Workshop on Number Theory and Cryptography --- Open Problems} (2006), attachment 6240, PDF page 38: \url{https://www.ipam.ucla.edu/programs/workshops/workshop-i-number-theory-and-cryptography-open-problems/?tab=schedule}. Registry access date 2026-07-29; PDF and legacy-PPT visual audit. Provenance reports 19 attachments, 17 PDFs and two legacy PPT decks, with repetitions and solved definitions excluded. Present status and redistribution terms require review.
\item The supplied summary cites Micciancio (2002) for compact knapsack/cyclic-lattice constructions and Lyubashevsky--Peikert--Regev (2010) for ring-LWE, without full bibliographic entries or precise assumptions and parameter bounds. These reports were not independently checked.
\end{itemize}
\end{document}
